\section{Bit complexity and the improvement over global branching}
\label{sec:complexity}

Let $V$ be the number of tree vertices, $b_v=|I_v|$ the number of original
coordinates at $v$, $\delta_v$ its degree, and $k_v=b_v+\delta_v$. For a fixed
incident sign assignment $\sigma$, let $M_{v,\sigma}$ be its local equation
matrix and put
\[
 d^{\rm raw}_{v,\sigma}=\dim\ker M_{v,\sigma},\qquad
 d_{v,\sigma}=\dim\operatorname{span}
       \{z\geq0:M_{v,\sigma}z=0\}.
\]
These need not coincide: forced-zero coordinates can lower the intrinsic
feasible dimension. Let $N$ denote the total binary encoding length of the
source, supplied decomposition, and objective; using the explicitly compiled
model instead changes this length by at most a polynomial factor.

\subsection{The delivered enumeration bound}

\begin{lemma}[Active-zero-set ray enumeration]
\label{lem:enum}
For a rational matrix $M$ with $k$ columns and kernel dimension $d>0$, every
extreme ray of $\{z\geq0:Mz=0\}$ is found by selecting $d-1$ coordinate
functionals that are independent on $\ker M$, setting them to zero, and solving
the resulting one-dimensional kernel. Testing all $\binom{k}{d-1}$ selections,
checking signs and forbidden pairs, and removing duplicates is complete.
\end{lemma}

\begin{proof}
At an extreme ray $r$, the face obtained by setting every coordinate outside
$\supp r$ to zero has dimension one. On $\ker M$, those zero-coordinate
functionals therefore have rank $d-1$. Select a basis of them. Their common
kernel in $\ker M$ is exactly the ray's linear span, so one of the two
orientations of the computed direction is $r$. Conversely, a nonzero
nonnegative direction whose support-restricted kernel is one-dimensional
spans an extreme ray. If $d=1$, the empty selection is used. If $d=0$, the cone
is zero and there are no rays.
\end{proof}

Define the explicit local enumeration budget
\begin{equation}\label{eq:Xi}
 \Xi=\sum_v\sum_{\sigma\in\{-1,1\}^{\delta_v}}
 \left(1+\mathbf1_{\{d^{\rm raw}_{v,\sigma}>0\}}
 \binom{k_v}{d^{\rm raw}_{v,\sigma}-1}\right).
\end{equation}
We use the convention $\binom{k}{-1}=0$. The leading $1$ charges the kernel
computation even for an empty cone.
The delivered code implements Lemma~\ref{lem:enum}. It does not implement a
polynomial-time LP engine or complete-support LP preprocessing.

\begin{theorem}[Implemented exact-query bound]
\label{thm:implemented}
Compilation, local enumeration, scalar elimination, and construction of one
positive witness have bit complexity
\[
 \poly(N)\,\Xi.
\]
In particular, with $h=\max_v(b_v+2\delta_v)$, the cost is bounded by
$\poly(N)2^h$. Thus a polynomial-size supplied source decomposition with
$h=O(\log N)$ gives a polynomial query algorithm, and
$h=O((\log N)^2)$ gives $N^{O(\log N)}$.
After local preparation, further objectives and original coordinate deletions
cost polynomial time in $N$ and the number of retained local rays.
\end{theorem}

\begin{proof}
There are $2^{\delta_v}$ local chambers at $v$. Each zero-set candidate uses
polynomially many exact rational arithmetic operations on matrices of
polynomial encoding length. Standard elimination and determinant bounds give
polynomial bit length for every local kernel direction and primitive ray.
Each candidate is either discarded or adds one ray, so all list sizes are
charged to \eqref{eq:Xi}. The dynamic recurrence scans those lists a constant
number of times per incident edge.

It remains to justify that rational message arithmetic does not escape the
bound. The following lemma supplies a polynomial bound on message and witness
bits. Finally, $\binom{k}{j}\leq2^k$, so summing over signs gives at most a
polynomial factor times $2^{b_v+2\delta_v}$ at each node. Repeated queries only
filter and rescore the already complete local lists.
\end{proof}

\begin{lemma}[No hidden expansion in message arithmetic]
\label{lem:bits}
Every finite message and every selected normalized witness has a rational
encoding of polynomial length in $N$. The scalar algorithm need not enumerate
normal discs or global compatibility histories to construct these numbers.
\end{lemma}

\begin{proof}
One proof uses the normalized subtree problem from
Theorem~\ref{thm:message}. After fixing its compatibility and scalar signs, it
is a rational standard-form system of polynomial dimension with one parent
coordinate set to $1$ in absolute value. Whenever its finite optimum is
attained, it is attained at a basic feasible solution. Clearing row denominators
and applying Cramer's rule bounds numerator and denominator bits by a polynomial
in the system's total encoding. The source compiler already has polynomial-size
coefficients. The objective coefficients on original coordinates are part of
the input; child slopes need not be included in this global subtree argument.
Thus every exact optimum value has polynomial encoding length.

For the actual selected recipe, each local coordinate-to-parent ratio is a
ratio of polynomial-size primitive-ray entries. A reconstructed coordinate is
a product of such ratios along a simple tree path, times one origin-ray entry.
There are at most $V$ factors. Numerator and denominator bit lengths add under
multiplication, so they remain polynomial. The objective is a sum of at most
$n$ products with the original rational coefficients. Reduced fractions in the
implementation prevent irrelevant common factors from accumulating.
\end{proof}

This is a bit bound, not a unit-cost arithmetic argument. A gain such as
$3^{1024}$ occupies only $1624$ bits, while the associated number of admissible
branch histories may be $2^{1024}$. Neither quantity is used as an iteration
bound in the scalar recurrence.

\subsection{A sharper local intrinsic-dimension theorem}

The following refinement uses the complete-support and dimension-drop principles
already available in report 87~\cite{repo87}. It is useful as a target for a
future LP-backed implementation.

\begin{lemma}[Compatibility cover by intrinsic dimension]
\label{lem:dimensioncover}
A nonnegative rational cone of intrinsic dimension $d$ with forbidden coordinate
pairs admits a complete binary compatibility search with at most
$2^{d+1}-1$ nodes, using polynomial-time rational LP at each node. On each
compatible leaf, feasibility or linear optimization under a specified
normalization is an ordinary LP problem.
\end{lemma}

\begin{proof}
At a node, compute which coordinates can be positive in the current cone.
This can be done by one feasibility query per coordinate, or by the complete
primal--dual support certificate of report 87. If no forbidden pair has both
coordinates live, every point of the cone is compatible. Otherwise choose a
forbidden pair $i,j$ with both coordinates live. Adding their positive witnesses
shows that they can be simultaneously positive in the convex cone.

Branch on $x_i=0$ or $x_j=0$. Every compatible point belongs to one child.
Each coordinate functional is nonzero on the current feasible span; hence each
child has intrinsic dimension at least one smaller. The depth is at most $d$,
and the node bound follows. All nodes use the original rational equations plus
coordinate zeros. Polynomial-time rational LP~\cite{lp} gives polynomial bit
cost per node. Normalized optimizations are performed only after a compatible
leaf is obtained, so no objective-dependent branching assumption is needed.
\end{proof}

\begin{theorem}[Sum of local intrinsic-dimension costs]
\label{thm:intrinsic}
The scalar-tree positive-objective problem can be decided in bit time
\[
 \poly(N)\sum_v\sum_{\sigma\in\{-1,1\}^{\delta_v}}
                  2^{d_{v,\sigma}}.
\]
This is a sum over local searches, not the product of their sizes. With
$D=\max_{v,\sigma}d_{v,\sigma}$ and $\Delta=\max_v\delta_v$, the bound is
$\poly(N)2^{D+\Delta}$. In particular, $D+\Delta=O((\log N)^2)$ is a
quasi-polynomial criterion even when the global matching-cone dimension is
linear in $N$.
\end{theorem}

\begin{proof}
At a node and sign chamber, unavailable child directions and disabled source
coordinates give coordinate faces and cannot increase intrinsic dimension.
Use Lemma~\ref{lem:dimensioncover} to cover the compatible local cone. On each
leaf, test for a positive zero-parent-port direction and optimize the effective
objective with parent absolute value one when feasible. If an optimum is
unbounded, its recession direction supplies the already decisive zero-port
positive witness. Theorem~\ref{thm:message} then computes the exact parent
message from these local outcomes.

Each subtree is processed once, and its compatibility history is replaced by
its message before the parent is solved. There is no enumeration of tuples of
leaf choices from different subtrees. Lemma~\ref{lem:bits} bounds the intermediate
coefficients, giving the displayed bit cost. The final estimate sums at most
$V2^\Delta$ chamber costs bounded by $2^D$.
\end{proof}

To retain the connected-witness conclusion with an LP backend, extract a vertex
of an optimal normalized local face rather than averaging tied optima. A
positive-port vertex corresponds to a cone ray; a positive zero-port search
can use a sum-of-coordinates normalization. Vertex extraction inside a rational
optimal face uses polynomially many exact LP or linear-algebra operations.
The delivered ray backend already chooses a single ray, including when several
rays have the same score.

\subsection{A strict separation from a global search budget}

Consider $m$ stages with nonnegative variables $t_0,\ldots,t_m,u_i,v_i$ and
\begin{equation}\label{eq:gainchain}
 t_{i-1}=u_i+v_i,\qquad t_i=2u_i+3v_i,\qquad u_iv_i=0
 \quad(1\leq i\leq m).
\end{equation}
At each stage, the only admissible nonzero choices transmit a gain of $2$ or
$3$. Consecutive stages share one nonnegative scalar. Eliminating the intermediate $t_i$ into scalar interfaces, and retaining $t_0$
and $t_m$ as endpoint coordinates, gives the disjoint-coordinate model used in
the code. The eliminated values are uniquely recovered and are automatically
nonnegative. Their objective coefficients can be absorbed into the adjacent
stage equations. The model has two endpoint rays and two admissible rays per
stage.

\begin{proposition}[Exponential global ray family, linear local inventory]
The unrestricted matching cone in \eqref{eq:gainchain} has dimension $m+1$.
Its admissible extreme rays are in bijection with the $2^m$ binary gain
histories. Nevertheless, every linear objective on its original variables can
be tested for positivity by the scalar recurrence with $2m+2$ retained local
rays and polynomial bit cost.
\end{proposition}

\begin{proof}
There are $3m+1$ variables and $2m$ independent equations. Independence follows
successively from the variables $t_m,t_{m-1},\ldots$ and the stage equations,
or by parameterizing the linear solution space using $t_0$ and $v_1,\ldots,v_m$.
The cone has a strictly positive point, obtained by taking $u_i=v_i>0$ at every
stage, so its intrinsic dimension is the same $m+1$.

For a nonzero compatible vector, $t_0>0$: if $t_0=0$, nonnegativity forces
$u_1=v_1=0$, hence $t_1=0$, and induction makes the entire vector zero. Normalize
$t_0=1$. At each stage exactly one of $u_i,v_i$ is positive, yielding its binary
history. Fixing that history makes the matching solution unique up to scale,
so the corresponding coordinate face is a ray. Different histories have
different original supports even when their final gains happen to agree.
The local ray counts and the scalar algorithm now prove the final assertion.
\end{proof}

For the objective $t_m-3^m t_0$, the answer is nonpositive; changing the
coefficient to $3^m-1$ gives a positive witness. The family provides a strict
comparison with a flat complete enumeration of all compatibility histories.
It is also analytically soluble by the elementary maximum-gain formula. It is
not offered as a difficult knot family, a lower bound on every algorithm, or a
novel closed-form optimization problem.

\section{A sharp warning for two-coordinate interfaces}
\label{sec:twoport}

The success of one scalar does not justify summarizing a two-dimensional
interface by a fixed number of slopes. The obstruction already occurs in a
constant-width binary counter.

\begin{theorem}[Exponential literal cone messages in interface dimension two]
\label{thm:twoport}
For every $m\geq1$, there is a homogeneous rational system with $O(m)$ variables
and constraints, coefficients bounded by $2$ in absolute value, and local
forbidden pairs, whose interface projection onto $(s,w_m)$ is exactly
\[
 P_m=\bigcup_{k=0}^{2^m-1}\{(s,ks):s\geq0\}.
\]
Any representation of $P_m$ as a finite union of convex cones requires at least
$2^m$ nonzero cones. The $2^m$ directions are distinguishable by homogeneous
linear contexts with a positive-objective query.
\end{theorem}

\begin{proof}
Take $s,u_i,v_i,w_i\geq0$, fix $w_0=0$, and impose
\[
 u_i+v_i=s,\qquad u_iv_i=0,\qquad w_i=2w_{i-1}+v_i
 \quad(1\leq i\leq m).
\]
When $s>0$, each $v_i$ is either $0$ or $s$. Thus $w_m/s$ is an arbitrary
$m$-bit integer, and every integer between $0$ and $2^m-1$ occurs. When $s=0$,
all coordinates vanish. This proves the projection formula. Consecutive
counter stages communicate only the pair $(s,w_i)$, so the phenomenon persists
with two-dimensional separators and bounded local equations.

If a convex cone contained nonzero vectors on two different displayed rays,
it would contain their positive combinations. Their slopes range over an
interval, including slopes not in the finite set of integers. Therefore a
convex cone contained in $P_m$ meets at most one nonzero displayed ray. At least
$2^m$ cones are required.

For a specified integer $k$, attach the homogeneous equation $w_m=ks$ and
ask whether the objective $s$ is positive. This context accepts exactly the
$k$th direction. Its coefficient has only $O(m)$ bits.
\end{proof}

The lower bound is explicitly about a \emph{literal union of convex cones}.
The counter itself is a short symbolic description, so the theorem cannot be
used to rule out arithmetic circuits, automata, semilinear descriptions, or
specialized exact optimization in rank two. It pinpoints a representation
barrier, not an unrestricted complexity lower bound.

\section{Connection to complete unknot recognition}
\label{sec:geometry}

\subsection{The standard-coordinate interface}

For a validated triangulation with $t$ tetrahedra, the standard coordinates
have four triangle and three quadrilateral entries per tetrahedron. The three
forbidden pairs among each quadrilateral triple enforce admissibility. Matching
equations, a source-derived rational Euler functional, and an original triangle
coordinate set to zero fit our homogeneous problem directly.

The classical positive-Euler/crushing recognition framework searches a
one-vertex knot-exterior triangulation for a connected positive-Euler normal
surface other than the vertex link; an essential-boundary disc proves
unknottedness, while the other permitted outcomes are handled by the geometric
reduction~\cite{bo}. Our source-bound ray theorem supplies connected witnesses
for that interface. It does not replace the essential-boundary test or the
proof that the triangulation is the exterior of the input knot.

In a one-vertex triangulation, the vertex-link vector has every triangle
coordinate equal to one. A connected normal surface different from the vertex
link has some triangle coordinate zero: otherwise a vertex-link copy can be
removed near the vertex as a disjoint normal component. Thus the search can
use the disjunction of at most $4t$ zero-triangle queries. Primitivity and
connectedness are important here; merely producing a positive Euler sum would
allow an irrelevant vertex-link component to dominate the objective.

For a genuine knot exterior, an embedded connected positive-Euler surface is
a sphere or a disc: an embedded projective plane does not occur in the
3-sphere. A sphere, a boundary-parallel disc, and an essential-boundary disc
must still be distinguished by the geometric layer. The generic algebraic
status \pos\ is deliberately not named \textsf{UNKNOT}.

\subsection{An explicit conditional running-time criterion}

\begin{theorem}[Width-controlled recognition criterion]
\label{thm:recognition}
Consider a correct positive-Euler/crushing recognition implementation operating
on a diagram of binary input size $N_0$. Assume the following contracts.
\begin{enumerate}[label=(\roman*),leftmargin=2em]
\item Source construction, one-vertex conversion, geometric checking, and
reduction have total polynomial bit cost outside their positive-Euler queries.
The number and encoding sizes of query sources are polynomial in $N_0$.
\item A decomposition producer, including all failed proposals, has total cost
$N_0^{O(\log N_0)}$ and returns source-verified, conflict-local rank-one tree
models for every query, or equivalent coordinate restrictions with authenticated
extension back to the source.
\item Across these query models, the implemented budget satisfies
$\sum\Xi\leq N_0^{O(\log N_0)}$.
\end{enumerate}
Then replacing the positive-Euler search by the delivered scalar ray method
preserves complete recognition and gives total running time
$N_0^{O(\log N_0)}$. A sufficient local condition for (iii) is polynomially many
bags with $b_v+2\delta_v=O((\log N_0)^2)$. For a polynomial-time LP implementation,
the alternative sufficient condition is
$d_{v,\sigma}+\delta_v=O((\log N_0)^2)$.
\end{theorem}

\begin{proof}
Source compilation preserves every query. Scalar elimination is complete in
both directions, and selected primitive rays give connected source surfaces.
The recognition implementation therefore receives exactly the required kind
of answer and continues through its already correct geometric steps. The
query costs sum by Theorems~\ref{thm:implemented} and \ref{thm:intrinsic}; the
producer cost and all nonquery costs are included explicitly. The displayed
sufficient conditions follow by substituting the local bounds.
\end{proof}

A single bag containing every source column always gives an exact model with
no interface condition. This supplies a complete exponential fallback at the
algebraic query level. It does not give the width assumption of the theorem.
Neither the one-vertex condition, normal matching, nor a small triangulation
dual-treewidth by itself has been proved here to imply that assumption.

The classical strictly decreasing tetrahedron count gives a particularly clear
special case: at most $t$ reduction stages and at most $4t$ triangle anchors
per stage give $O(t^2)$ anchored queries. Prepared local rays can be reused
across anchors on a fixed source. Reuse across a changed source is not covered.

\subsection{Relation to Lackenby's hierarchy program}

The scalar theorem is a search primitive, not a substitute for the entire
hierarchy. To use it for later hierarchy pieces, the source constraints must
still express the correct surfaces and inherited boundary conditions. A
marked-boundary objective or charge may introduce several interface quantities,
nonhomogeneous conditions, or cross-bag restrictions. Such changes must be
proved compatible with the message contract rather than inferred from a small
raw-cut encoding.

In particular, a proof of short hierarchy depth does not itself prove small
scalar width, and small scalar width at one stage does not itself control
hierarchy resets. The most useful next connection would be a source-aware
normal form for the hierarchy pieces in which most of the matching system
splits by rank-one cuts and only a polylogarithmic residual core remains.
That is a concrete additional theorem to seek, not a conclusion established
by the present experiments.
